\documentclass[12pt]{article}
\usepackage{amsmath,amssymb,amsfonts}
\usepackage[margin=1in]{geometry}

\begin{document}

\textbf{Chapter 6, Problem 1.} Let
\[
f(z) = \frac{x}{x^2 + y^2} \quad \text{for } x \neq 0, \quad f(0) = 0.
\]
Determine where, if anywhere, this function is (a) differentiable, (b) analytic.

\bigskip
\textbf{Solution.}

We write $z = x + iy$ and identify $f(z) = u(x,y) + iv(x,y)$. Here $f$ is real-valued, so $v(x,y) = 0$ and
\[
u(x,y) = \frac{x}{x^2 + y^2}.
\]

For $f$ to be differentiable at a point, the Cauchy--Riemann equations must hold:
\[
\frac{\partial u}{\partial x} = \frac{\partial v}{\partial y}, \qquad \frac{\partial u}{\partial y} = -\frac{\partial v}{\partial x}.
\]

Since $v = 0$ everywhere, we need $\partial u/\partial x = 0$ and $\partial u/\partial y = 0$.

Compute:
\[
\frac{\partial u}{\partial x} = \frac{(x^2+y^2) - x \cdot 2x}{(x^2+y^2)^2} = \frac{y^2 - x^2}{(x^2+y^2)^2},
\]
\[
\frac{\partial u}{\partial y} = \frac{-2xy}{(x^2+y^2)^2}.
\]

Setting $\partial u/\partial y = 0$ gives $xy = 0$, so $x = 0$ or $y = 0$.

Setting $\partial u/\partial x = 0$ gives $y^2 = x^2$, so $y = \pm x$.

Both conditions together: $xy = 0$ and $y = \pm x$. From $xy = 0$ we have $x = 0$ or $y = 0$. If $x = 0$, then $y = 0$; if $y = 0$, then $x = 0$. So the only candidate point is $x = y = 0$, but $u$ is not even defined at the origin in the usual sense (and $f(0)$ is separately defined as 0).

At the origin, we check differentiability directly. Approaching along the real axis ($y = 0$):
\[
\frac{f(h) - f(0)}{h} = \frac{1/h}{h} = \frac{1}{h^2} \quad (\text{for real } h \neq 0),
\]
wait -- for real $h$, $f(h) = h/(h^2) = 1/h$, so $[f(h)-f(0)]/h = 1/h^2 \to \infty$. The derivative does not exist at the origin.

Therefore, $f(z)$ is \textbf{nowhere differentiable} (in the complex sense), and consequently \textbf{nowhere analytic}.

\end{document}
